% =============================================================================
%  CHAPTER 8 -- CONCLUSION                                            [v4 OWNS]
% -----------------------------------------------------------------------------
%  v4.0 (2026-09-24): rewritten by v4 from the author's general conclusion and the
%  thesis-in-three-steps (../GUIDE_20260924_results_storyline_author.md, v3.96/v3.97) on
%  Chapters 5-6 at v3.98 and Chapters 1-4 at v2.27 (via v3). v3's concise 20-09 draft
%  (FROM_v3_v3.98_20260924_201920/chapters/08_conclusion.tex) was stale against Ch 6 and is
%  not reused. Labels sec:conc:summary, sec:conc:rqs, sec:conc:future are kept;
%  sec:conc:practice is new. Every number restates one of Chapter 6.
%  Individual changelog: changelogs/v4.0_20260924_init_ch7_ch8_appendix.md
% =============================================================================

\chapter{Conclusion}\label{ch:conclusion}

\section{Summary}\label{sec:conc:summary}

This thesis replaced the diffusion model of \ac{DPCC} with three flow-based objectives,
instantaneous-velocity matching, analytic average-velocity matching and consistency-interpolated
average-velocity matching, and set endpoint projection beside the per-step projection of \ac{DPCC},
keeping the temporal U-Net, the projection program, the data and the evaluation of the baseline. The
step budget thereby moved from training to inference. The comparison was made in three steps: on the
obstacle-avoidance benchmark of \ac{DPCC} with its setup unchanged, on a vision-conditioned alignment
task built for this thesis, and on a quadrotor benchmark of three scenes built for it as well.

On the benchmark of \ac{DPCC} the average-velocity models Pareto-dominate its diffusion model with its
own U-Net: the same success with constraint satisfaction, fewer control steps, and about a thirtieth
of the time per action, twelve times less than a baseline retrained at two denoising steps.
Instantaneous-velocity matching dominates the baseline as well, so the saving belongs first to a
flow-based sampler at a small budget and then to the average-velocity objective, which adds the fewest
control steps. On D3IL-aligning the dominance repeats where the task separates the models: analytic
average-velocity matching is the only model that moves the box in every context, it closes most of
the distance to the target, and it is the only objective that improves with the budget, while the
consistency-interpolated model, instantaneous-velocity matching and the baseline barely move the box.
Endpoint projection has no step to guide at the one-evaluation operating point of D3IL-avoiding and
nothing to win there; on D3IL-aligning, run at twenty evaluations, it is the better projector on the
constraint, keeping every context violation-free at the price of per-step projection.

The quadrotor scenes carry the framework to a plant that holds nothing by itself. UAV-pillars flies
the operating points of D3IL-avoiding and keeps their result on the declared constraints, and shows
what the plant changes: the vehicle tracks more closely than the arm, and a collision with a pillar the
constraints leave out ends the flight. UAV-corridor shows that the projection scheme of \ac{DPCC}
carries over to three dimensions, against constraints that leave the horizontal plane and that no
demonstration satisfies, and that on straight-line demonstrations the generative model does not
matter, the budget does; there the two projectors trade against each other without one dominating.
UAV-s-curve shows that the controller matters: the same plans succeed or fail by the controller that
flies them, and the cascaded geometric controller used throughout is the better of the two, at a
twenty-fourth of the cost.

With human demonstrations the average-velocity models lead, on the harder tasks the
model and the projection matter together and endpoint projection helps where the task needs a large
budget; with simple generated demonstrations instantaneous-velocity matching is enough and the
objectives are equivalent; and on the quadrotor the plant and its controller matter as much as the
planner.

\section{Answers to the Research Questions}\label{sec:conc:rqs}

\begin{description}
  \item[RQ1 --- Generative model.] Yes, with a margin. With the backbone, the data and the evaluation
    held fixed, the flow-based models reach the diffusion baseline's success with constraint
    satisfaction at one network evaluation, where the diffusion model, run away from its trained
    budget, loses it: on D3IL-avoiding consistency-interpolated average-velocity matching holds
    $1.000$ in $59.2$ control steps at $18.1$\,ms per action against the baseline's $1.000$ in $70.1$
    at $553.4$, and instantaneous-velocity matching $1.000$ in $67.0$ at $17.3$. Among the objectives
    the order, wherever the task separates them, is the average-velocity models, then
    instantaneous-velocity matching, then the diffusion model. On D3IL-aligning analytic
    average-velocity matching closes $84\,\%$ of the distance to the target against $14$, $10$ and
    $-2\,\%$ for the others, and under projection keeps nine contexts violation-free against the
    baseline's four at a third of its time. Instantaneous-velocity matching and the diffusion model
    are level on outcome there and on D3IL-avoiding, and differ in cost. Between the two
    average-velocity objectives the consistency-interpolated one is ahead by one episode in thirty on
    D3IL-avoiding and behind by most of the task on D3IL-aligning, where it does not improve with
    the budget; the analytic objective leads or holds in every environment. Where the demonstrations
    are simple, on UAV-corridor and UAV-s-curve, the objectives are equivalent or
    instantaneous-velocity matching leads, and the budget orders the configurations instead. One
    evaluation saturates D3IL-avoiding, D3IL-aligning needs twenty, and on the s-curve more
    evaluations lose flights.
  \item[RQ2 --- Constraints.] Endpoint projection guides the sample only where the budget and the
    threshold leave a sampling step before the terminal one, $n\sidx{guide} = \lceil \eta\nfe \rceil
    - 1$: at one evaluation never, at two at the default threshold not, and there it computes the same
    plan as per-step projection at $2.45$ times the cost. Where it has steps to guide it is the better
    projector on the constraint at a comparable price: on D3IL-aligning at twenty evaluations it keeps
    all ten contexts violation-free where per-step projection keeps nine, in seven of nine matched
    comparisons, and at ten evaluations at $37\,\%$ of per-step projection's time; on D3IL-avoiding at
    three evaluations it halves the time per control step of the two average-velocity models at
    $1.000$. It does not reach the goal in fewer control steps, and on UAV-corridor it trades: fewer
    and shallower violations at twenty evaluations, cheaper over the hump with one candidate and
    dearer under the tilt, on longer flights. It helps on the hard task that needs the budget,
    D3IL-aligning, and not on D3IL-avoiding, which is solved at a budget where it has no step to act.
  \item[RQ3 --- Transfer.] The answers hold where the demonstrations are of the same kind, and the
    plant and the demonstrations change the rest. Analytic average-velocity matching leads again on
    camera observations and contact, and the declared constraints transfer exactly from the
    manipulator to the quadrotor on UAV-pillars, where the two average-velocity models keep their
    order. The quadrotor adds what the manipulator could not show: an undeclared obstacle ends a
    flight, the controller decides whether a plan on a route of sharp turns is flown at all, and on
    demonstrations that carry one path forward from every position the objectives cannot be told
    apart. The projection scheme itself carries over to three dimensions and to constraints that leave
    the horizontal plane.
\end{description}
\dataref{Every number restates \autoref{sec:res:summary}: \autoref{tab:avoiding-conclusion},
\autoref{tab:va-models}, \autoref{tab:va-projection-models}, \autoref{tab:avoiding-projectors},
\autoref{sec:res:uav:projection}, \autoref{tab:uav-pillars-raw}, \autoref{tab:uav-controller}.}

\section{Towards Deployment}\label{sec:conc:practice}

Three properties of the framework carry into practice as it stands. The constraints enter only at
sampling time, through the projection; nothing about them reaches training, so a new restriction is a
change of the constraint set and not a training run, and the feasible set can be replaced between two
replans, as the switched walls of the quadrotor scenes are (\autoref{sec:method:env:uav}). The step
budget is a deployment choice: the same checkpoint of a flow-based model is read at one evaluation on
D3IL-avoiding and at twenty on D3IL-aligning, where a diffusion model needs one checkpoint per budget.
And the plan is a sequence of positions that a tracking controller follows, which is what let one
planner fly the quadrotor unchanged.

Two quantities bound that use. A plan is held inside the declared geometry only, so the constraint set
has to name every obstacle that matters, with a margin that covers the vehicle's extent, and the
channel it leaves has to be flyable by the controller: on the corridor the commanded position runs
$0.49$ to $0.61$\,m ahead of the vehicle, on the s-curve about $0.3$\,m, and the vehicle follows the
commanded path to within about a centimetre sideways. And the time to compute one action, $18$\,ms at
the operating point of D3IL-avoiding and $266$ to $275$\,ms on D3IL-aligning on the node of record,
prices the planner alone; the execution adds its own, $5$\,ms per control step for the cascaded
geometric controller and $125$\,ms for MuJoCo MPC on the s-curve, and the manipulator's inverse
kinematics and joint controller were not timed. \guard{No loop of this thesis was run in real time;
the simulation waits for the planner (\autoref{sec:method:geometric}).}

\section{Future Work}\label{sec:conc:future}

\begin{description}
  \item[Demonstrations for the quadrotor.] The demonstrations of UAV-corridor and UAV-s-curve are
    generated, one path forward from every position, and the s-curve's plans cut the corner that the
    demonstrated route clears. Human demonstrations of the quadrotor scenes, from a pilot in the loop
    of a simulator or from flight logs, would give the average-velocity objectives the multimodal data
    on which they separate and the s-curve a route its plans do not cut; a simulator with a rendered
    scene and a pilot interface is the way to record them without a vehicle.
    \hole{Author: name the simulator (NVIDIA Omniverse / Isaac Sim, as in notes.txt) or keep it generic.}
  \item[The controller.] On the s-curve the projector binds the measured position while the
    corridor's binds the commanded one; the projector that binds both was not flown there. A plan that
    carries a velocity setpoint, or a controller that reads the next positions of the plan as
    feed-forward, would close the lead between the commanded and the flown position from which the
    hand-over residue and the corner cut come.
  \item[Real-time execution.] The execution stack has to be timed and the loop run without waiting for
    the planner: the planner at one evaluation, the projector solved in parallel over the candidates,
    or with a single candidate where the rule selects the same plan, and the controller at its own rate.
  \item[Constraints from perception.] The projector consumes a short list of primitives with metric
    extent, halfspaces, keep-out disks and spheres and a workspace box, each in the planner's frame and
    each with its feasible side. A perception front-end that emits that list from depth or range
    sensing, and a map that keeps it beyond the field of view, is what a deployment outside the
    simulator needs; the uncertainty of the detector belongs in the inflation margin the projector
    already carries. The first experiment needs no detector: perturb the geometry the projector
    receives, by an offset, a radius error, a missing or a spurious primitive, and record where success
    with constraint satisfaction breaks.
  \item[Re-tasking at deployment.] Because the constraints never enter training, the same checkpoint
    can be asked to satisfy a new set: a closed passage, a keep-out volume, a tighter clearance. What
    has to be measured is the envelope, the slack a new set may leave before the projector pulls the
    plan off the data and success collapses rather than degrades, and the transfer of one checkpoint to
    a geometry its demonstrations never contained.
  \item[The alignment target.] The optional terminal cost of endpoint projection
    (\autoref{sec:method:hardflow}), zero in this thesis, is where the target of D3IL-aligning could
    enter the projection; constraint satisfaction would then stop costing completion.
  \item[The sampler of instantaneous-velocity matching.] Sampling it from the scale it was trained at
    needs no retraining and would settle its order against the average-velocity models on
    D3IL-aligning and the quadrotor scenes.
  \item[Seeds.] More training seeds on D3IL-aligning and the quadrotor scenes.
\end{description}
\srcnote{Outlook material: ../future\_work/OUTLOOK\_20260910\_perception\_to\_constraints.md (the
projector's interface, the perturbation experiment) and
OUTLOOK\_20260910\_zero\_shot\_constraint\_retasking.md (constraints never enter training, verified in
code; the envelope experiment); the author's v4/notes.txt (expert data for the quadrotor, the
sim-to-real timing question, the selection-rule discussion). Nothing here is a result.}
